\chapter{Introduction}
\section{Objective and structure of the thesis}
The objective of this thesis is ...
\begin{itemize}
    \item In the first chapter, the main properties of Dark Matter are presented, together with its most relevant candidates: WIMPs. Then the most significant astrophysical evidence is discussed, on scales ranging from galaxies and galaxy clusters to the entire Universe. In addition, theories of modified gravity are briefly introduced, developed as an alternative to the Dark Matter hypothesis. Finally, the theoretical framework for DM indirect searches is illustrated and the main astrophysical environments investigated are described.
          The primary source of information for this chapter is \emph{Dark Matter} \cite{cirelli2025darkmatter}, one of the most complete review on Dark Matter written to date.
    \item In the second chapter,
\end{itemize}

All the code developed and used in this work for numerical calculations is available on \href{https://github.com/IvanDalena/Bachelor-thesis}{GitHub}.
 
Cose da sistemare:


-Riproduci tutti i dati tramite codice fatto da te, in particolare i J factor della tab nel GC






\chapter{Dark Matter}

Evidence for the existence of Dark Matter comes from a wide range of astronomical scales: from a few kiloparsecs, corresponding to the typical size of a small galaxy, up to nearly the entire Universe.

In this chapter, after a brief description of the properties of Dark Matter and its leading candidates, this evidence is presented over increasing scales: starting from the study of rotation curves of spiral galaxies, to the mass of galaxy clusters, and finally to the Cosmic Microwave Background (CMB) and Large Scale Structures. 

Since all this evidence can be explained through gravitational interactions, alternative theories of gravity have been developed: the most prominent one is Modified Newtonian Dynamics (MOND). However, none of these theories is currently able to explain all astrophysical observations.

The present cosmological Dark Matter density, averaged over the total matter-energy content, is estimated to be $\Omega_{DM} = 0.264 \pm 0.003$. Considering only the matter content, this value rises to about 84\%. In \cref{UniverseContent} is shown a comparison between the content of the Universe today and at the time of photon decoupling, the moment when photons decouple from matter and start to travel freely.



\begin{figure}[H]
    \centering
    % Parte sinistra: Porta OR
    \begin{minipage}{0.47\textwidth}
        \centering
        \includegraphics[width=\textwidth]{UniverseContentToday.png} 
        
    \end{minipage}
    \hfill % Crea spazio tra le due immagini
    % Parte destra: Porta AND
    \begin{minipage}{0.47\textwidth}
        \centering
        \includegraphics[width=\textwidth]{UniverseContentCMB.png} 
       
    \end{minipage}

    \vspace{0.3cm} % Spazio estetico tra foto e didascalia
    \caption{\emph{Content of the Universe today (left) and, for comparison, at the time of photon decoupling (right). The different components evolve differently in cosmology \cite{cirelli2025darkmatter,Planck2018}.}}
    \label{UniverseContent}
\end{figure}



\section{Properties of Dark Matter and WIMPs}

From astrophysical and cosmological observations it is possible to determine the general properties of Dark Matter: it turns out that Dark Matter behaves as a kind of \emph{cold, non-interacting, stable matter with adiabatic inhomogeneities}.

\begin{itemize}
    \item \textbf{Matter} stresses the fact that DM behaves as matter in the cosmological evolution (i.e. its density decreases as the inverse of the volume).
   
    \item \textbf{Cold} means that DM behaves in large part as a non-relativistic fluid at the time of matter-radiation equality (MReq). Therefore, DM particles move slowly (with respect to the speed of light) and can attract each other and cluster. If they moved with velocities comparable to \emph{c}, the clustering would not be effective and structures would not form and grow.
   
    \item \textbf{Non-interacting} (or equivalently \textbf{collision-less}) means that the interactions of DM with itself or between DM and ordinary matter are small enough to be neglected, except for gravitational interaction. For this reason it is called "Dark Matter": the absence of interactions with light makes DM \emph{dark}.
    
    \item A consequence of the non-interacting nature is that Dark Matter is \emph{dissipation-less}: unlike ordinary matter, DM cannot emit electromagnetic radiation and therefore cannot easily dissipate its energy and cool down. This is the main reason why baryonic matter and Dark Matter behaved differently during cosmological evolution.
    
    \item \textbf{Stable} means that DM is present since the early phases of the Universe (estimated to be $t_{\mathrm{Uni}} \simeq 13.8\,\mathrm{Gyr} = 4.35 \times 10^{17}\,\mathrm{s}$). Its lifetime has to be over 10 orders of magnitude longer than the lifetime of the Universe.
    
    \item \textbf{Adiabatic} means that, on cosmological scales, its composition is the same everywhere. Moreover, DM and ordinary matter has the same primordial density inhomogeneities, indicating that the mechanism responsible for their generation is the same.
    
\end{itemize}

While there is broad consensus about the hypothesis that DM is made of one or more new particles beyond the Standard Model (SM) of particle physics, 
no definitive experimental confirmation has been obtained so far. Among the particles that satisfy these conditions, the most studied and well-motivated candidates are the \emph{Weakly Interacting Massive Particles} (WIMPs). 
WIMPs are assumed to interact through the weak force and gravity, with masses in the range $\mathcal{O}(1\text{--}100)~\text{GeV}$ and a cross section characteristic of the electroweak scale, which justifies their name.
It is assumed that in the early Universe WIMPs were in thermal equilibrium with the cosmic plasma, with an interaction rate $\Gamma$ much larger than the Hubble expansion rate $H$. When $\Gamma \sim H$, they decoupled from the primordial plasma, the so-called \emph{freeze-out}, breaking this equilibrium and fixing the average density of WIMPs in the Universe.
The \emph{WIMPs miracle} consists in the fact that, by requiring this relic density to reproduce the observed Dark Matter abundance, one obtains an annihilation cross section of $\langle\sigma v\rangle \simeq 3\times10^{-26}~\text{cm}^3/\text{s}$, compatible with electroweak-scale interactions.
No conclusive signal, despite an extensive search program that combines, often in a complementary way, direct, indirect, and collider probes, has been however detected so far. This situation might change in the near
future with the advent of even larger detection experiments \cite{profumo, Arcadi_2025}.

\section{Astrophysical evidence}

\subsection{Rotation curves of spiral galaxies}

In 1970, Vera Rubin and Kent Ford initiated a true revolution in the field of Dark Matter after analyzing the rotation curves of 
the Andromeda Galaxy (M31) \cite{Ford}.

It is known that spiral galaxies rotate around their vertical axis, therefore it is possible to infer their mass distribution from the rotation velocity, which can be measured experimentally through Doppler shift.
From Newtonian mechanics it is possible to derive the circular velocity $v_{\mathrm{circ}}$ of a body of mass $m$ at a distance $r$ from the center as a function of the mass $M$ enclosed within $r$, assuming spherical symmetry


\begin{equation} \label{circular}
\frac{m\,v_{\mathrm{circ}}^{2}(r)}{r}
=
\frac{G\,m\,M(r)}{r^{2}}
\quad \Rightarrow \quad
v_{\mathrm{circ}}(r)=\sqrt{\frac{G\,M(r)}{r}}\,.
\end{equation}

In spiral galaxies, most of the mass is concentrated relatively close to the center and along the disk arms. Therefore, beyond a certain radius $r$, $M(r)$ can be considered approximately constant and the circular velocity should follow a Keplerian decline $v_{\mathrm{circ}}(r) \propto r^{-1/2}$.
The crucial point is that the observations of a large sample of galaxies of this kind have shown that the rotation curves instead remain \emph{flat} out to large distances from the galactic centers. The results obtained by Rubin and Ford for the Andromeda galaxy are shown in \cref{RCM31}.


\begin{figure}[H]
\centering
\includegraphics[scale=0.45]{RCM31.png}
\caption{\emph{Rotation curve of the Andromeda Galaxy (M31) \cite{Ford}}.}
\label{RCM31}
\end{figure}

To explain this discrepancy, additional invisible mass must be present to prevent the peripheral stars from flying away and the galaxies from breaking apart. Assuming spherical symmetry, \cref{circular} shows that a constant velocity $v_{\mathrm{circ}}(r)$ is obtained assuming a matter \emph{halo} that has mass density $\rho(r) \propto \frac{1}{r^{2}}$ out to large $r$: 
in this way, $M(r)=4\pi \int dr'\,{r'}^{2}\rho(r') = 4\pi r$ so that the dependence of $v_{\mathrm{circ}}(r)$ on $r$ vanishes. Eventually, at even larger $r$, the halo is expected to die off and the rotation curve to start declining. However, no traces are typically available at such large distances \cite{cirelli2025darkmatter}.

To describe the Dark Matter distribution in more detail, numerical simulations are used: from $N$-body simulations, including only the Dark Matter component, one obtains a profile known as the Navarro-Frenk-White (NFW) \cite{NFW1996,NFW1997}


\begin{equation} \label{NFW}
\rho_{\mathrm{NFW}}(r)\propto \frac{1}{\dfrac{r}{r_s}\left(1+\dfrac{r}{r_s}\right)^2}\,
\end{equation}

where $r_s$ is a characteristic parameter called scale radius \cite{Lisanti:ICTP:2018:DM}.


%\begin{figure}[H]
%\centering
%\includegraphics[scale=0.45]{RC50Galaxies.png}
%\caption{\emph{Foto da \cite{Sofue_1999}}}
%\label{RC50Galaxies}
%\end{figure}

\subsection{Galaxy clusters}


The first evidence of a mass-to-light discrepancy dates back to 1933, when Fritz Zwicky studied the velocity dispersion of several galaxy clusters and obtained unexpected results \cite{Zwicky1933,Zwicky}.
Using the virial theorem, he estimated the mass of the clusters from the velocity dispersion of the galaxies, measured through Doppler shift, and found large discrepancies with the mass inferred from the observed luminous matter.
Zwicky is also credited with introducing the term \emph{Dark Matter}, although the concept itself had already been present in the scientific community for some years \cite{History}.

Particularly significant are the results obtained for the Coma cluster, for which Zwicky concluded that additional mass was required to gravitationally bind the galaxies.
It is possible to make these results quantitative: from the virial theorem, for a gravitationally bound system, the time averaged kinetic energy $\langle K \rangle$ and potential energy $\langle V \rangle$ are related by


\begin{equation}\label{eq:viriale}
\langle K \rangle = -\frac{1}{2}\langle V \rangle.
\end{equation}

Considering the one dimensional velocity dispersion of the galaxies, measured along the line of sight between the observer and the cluster, 
equal to $v_{\text{los}} \sim 10^{3}\ km/s$, and assuming an isotropic velocity distribution, one obtains the three dimensional dispersion


\begin{equation}
\langle v^{2}\rangle = 3\,v_{\text{los}}^{2}.
\end{equation}

Consequently, the time averaged kinetic energy can be expressed in terms of the total mass of the cluster $M$ as

\begin{equation} \label{Kinetic}
\langle K\rangle \;=\;\frac{1}{2}M\,\langle v^{2}\rangle.
\end{equation}

Modeling the cluster as a sphere of total mass $M$, uniformly distributed within a radius $R$, the gravitational potential energy is given by

\begin{equation}\label{eq:Vpot}
    \langle V\rangle = -\frac{3\,G M^{2}}{5\,R}.
\end{equation}

Substituting the expressions obtained for $\langle K \rangle$ and $\langle V \rangle$, respectively \cref{Kinetic} and \cref{eq:Vpot}, into \cref{eq:viriale} it is possible to derive an expression for the mass of the cluster

\begin{equation}
\frac{1}{2}M\langle v^{2}\rangle
=
-\frac{1}{2}\!\left(-\frac{3}{5}\frac{GM^{2}}{R}\right)
\qquad\Rightarrow\qquad
M=\frac{5R\,\langle v^{2}\rangle}{3G}
=\frac{5R\,v_{\text{los}}^{2}}{G}.
\label{eq:Mviriale}
\end{equation}

To determine the value of $M$, it is necessary to estimate the radius $R$, which can be obtained geometrically from the distance of the cluster from the Earth, $d \sim 100~\mathrm{Mpc}$, and its angular radius $\theta \simeq 1^\circ$. 
It is therefore possible to use the small-angle approximation

\begin{equation}
R \simeq d\,\theta \simeq 1.7\ \text{Mpc}.
\end{equation}

Substituting $v_{\mathrm{los}} \sim 10^{3}~\mathrm{km/s}$ and $R \simeq 1.7~\mathrm{Mpc}$ into \cref{eq:Mviriale}, one obtains a total mass of the Coma cluster of the order of

\begin{equation}
M  \simeq 2 \times 10^{15}\,M_{\odot}.
\end{equation}


It is crucial to compare this value with the baryonic mass of the cluster, which is of the order of $10^{14}\,M_{\odot}$, obtained by summing the contribution from the visible mass of the $\sim 10^3$ galaxies and the hot gas measured through $X$-ray observations.
It is thus observed that the mass responsible for the gravitational interaction is at least one order of magnitude larger than the estimated baryonic mass, implying the presence of additional, yet unknown, matter.

A further piece of evidence for the existence of DM comes from a system formed by two galaxy clusters in the process of colliding known as the \emph{Bullet cluster}, located 3.7 Gyr away.
As shown in \cref{BulletCluster}, $X$-ray emissions trace the regions where most of the baryonic matter is concentrated, which in this system is primarily in the form of hot gas. On the other hand, weak gravitational lensing allows the reconstruction of the total mass distribution of the system.
The comparison between these observations reveals a clear separation between baryonic matter and the mass responsible for the gravitational interaction, once again suggesting the presence of additional matter.
The special feature of the bullet cluster is that it shows the non-interacting nature of Dark Matter: it is supposed that in the past each of the two clusters had the visible matter and DM mixed together. The two objects collided 150 million years ago. Visible matter interacts significantly with itself, so that the hot gas from the two clusters experienced a collisional shock wave. DM, on the other hand, experienced negligible collisions with itself and with normal matter, such that it simply passed through.



% \begin{figure}[H]
% \centering
% \includegraphics[scale=0.5]{BulletClusterDX.png}
% \caption{\emph{Mappa a colori della distribuzione di gas caldo barionico, ricavata dall'emissione di raggi X. I contorni verdi mostrano la distribuzione della massa totale ottenuta tramite lensing gravitazionale debole. La barra bianca corrisponde a una lunghezza di 200 kpc \cite{Clowe_2006}.}}
% \label{BulletCluster}
% \end{figure}

\begin{figure}[H]
    \centering
    
    \begin{minipage}{0.47\textwidth}
        \centering
        \includegraphics[width=\textwidth]{BulletClusterSX.png} 
        
    \end{minipage}%
    \hfill %
    %
    \begin{minipage}{0.47\textwidth}
        \centering
        \includegraphics[width=\textwidth]{BulletClusterDX.png} 
       
    \end{minipage}%

    \vspace{0.3cm} 
    \caption{\emph{On the left, a color image of the merging clusters.
                   On the right, a color map of the distribution of hot baryonic gas, reconstructed from $X$-ray emission, is presented.
                   The green contours indicate the total mass distribution obtained through weak gravitational lensing. The white bar corresponds to a length of 200 kpc \cite{Clowe2006}.}}
    \label{BulletCluster}
\end{figure}

This kind of observations puts a severe strain on alternative interpretations where DM is replaced by a modification of gravity, because such modifications cannot get spatially separated from normal matter. In the next section, this these theories will nevertheless be briefly discussed.

In this context, it is also important to mention the Cosmic Microwave Background (CMB) and the Large Scale Structures, which provide further evidence in support of Dark Matter on cosmological scales, where the most convincing and precise evidence comes nowadays.
Dark Matter played a fundamental role as a catalyser in the formation of galaxies and galaxy clusters. Moreover, CMB temperature anisotropies would not look the way they do if it weren't for DM. This was accurately studied by the Planck collaboration, which measured the Dark Matter density from the analysis of the CMB power spectrum, obtaining $\Omega_c h^2 = 0.120 \pm 0.001$, in excellent agreement with the predictions of the $\Lambda$CDM model \cite{Planck2018}.



\subsection{Modified Newtonian Dynamics}

In February 1982, Mordehai Milgrom submitted three papers \cite{Milgrom1983a,Milgrom1983b,Milgrom1983c} which provided the
foundation for what would become the leading alternative to Dark Matter. 
Milgrom's proposal states that gravity (or inertia) does
not follow the predictions of Newtonian dynamics for accelerations smaller than the characteristic acceleration $a_0 = 1.2 \times 10^{-8}\,\mathrm{cm\,s^{-2}}$, as derived from galaxy rotation curves fits.
The MOND gravitational acceleration $a$ is related to the Newtonian acceleration $a_N$ through an interpolation function $\mu(a/a_0)$, such that
\begin{equation}
    a_N = a\,\mu\left(\frac{a}{a_0}\right)
\end{equation}
where $a_0$ is meant to be a new constant of physics.

By construction, the interpolation function has the asymptotic behavior
$\mu = 1$ for $a \gg a_0$, to recover the Newtonian expression in the strong field regime, and $\mu = a/a_0$ for $a \ll a_0$.
In this limit, the gravitational acceleration becomes
\begin{equation}
    a = \sqrt{a_N\,a_0} = \frac{\sqrt{GM a_0}}{r}
\end{equation}
with dependence $1/r$ instead of the Newtonian $1/r^2$ \cite{Scarpa_2006}.
Within this framework, MOND successfully reproduces the flat rotation curves of galaxies. However, the value $a_0$ is introduced as an empirical parameter adjusted to match observational data. In fact, a large variation in $a_0$ is found studying different samples of galaxies. 
This shows that MOND works well as a fitting procedure but not as a phenomenological
description of gravity on the scale of galaxies \cite{a0}.

Moreover, in its original formulation MOND suffered from several theoretical issues, such as the lack of energy conservation and the difficulty of a consistent relativistic extension.
Subsequently, new theories were developed, such as the AQUAdratic Lagrangian (AQUAL) \cite{AQUAL} and its relativistic extension RAQUAL, which are based on a Lagrangian formulation. Nevertheless, all these theories struggle to provide a unified description capable of simultaneously explaining the full range of experimental observations.



\section{Indirect detection}


Indirect detection experiments aim to probe the nature of Dark Matter by searching for an excess (above the presumed astrophysical contribution) in the flux of Standard Model particles produced through Dark Matter annihilation or decay. A schematic representation of the signal production mechanism is shown in \cref{DMannihilation}.
In this context, the most relevant messengers are photons (gamma rays and $X$-rays), neutrinos, and charged cosmic rays (such as positrons, antiprotons, antideuterons, and heavier antimatter nuclei).
For charged cosmic rays, particular attention is given to antimatter particles since they are less likely to be produced in astrophysical processes than the corresponding matter particles, leading to more favorable astrophysical background. However, their propagation through the Galaxy introduces additional systematic uncertainties, such as changes in their direction due to the galactic magnetic fields or energy losses while propagating.

First of all, it is necessary to assume that Dark Matter possesses at least one interaction channel producing Standard Model particles. While there is no experimental confirmation of this, it is a natural expectation in well-motivated scenarios such as the "thermal relic" one, where the relic abundance is directly related to its annihilation cross section \cite{Slatyer}.
In alternative scenarios, in which DM annihilates or decays exclusively into itself, conventional indirect detection would not be effective.

The flux of cosmic rays produced by Dark Matter primarily depends on the annihilation cross section $\langle\sigma v\rangle$ or, in the case of decay, on the particle lifetime $\tau$.
Another crucial factor is the spatial distribution of DM, which enters through line of sight integrals called \emph{J factor}: the resulting flux is therefore expected to be enhanced in regions of higher Dark Matter density, such as the Galactic Center or dwarf spheroidal galaxies, although these environments might be challenging to study. 



\begin{figure}[H]
\centering
\includegraphics[scale=0.6]{DMannihilation.png}
\caption{\emph{Annihilation of Dark Matter into Standard Model particles and generation of primary particle spectra \cite{cirelli2025darkmatter}}.}
\label{DMannihilation}
\end{figure}


\subsection{Gamma rays}

The gamma ray spectrum from DM annihilation can be classified according to their shape as follows
\begin{itemize}
    \item $\gamma$-ray lines, from the direct annihilation of two DM particles into $\gamma\gamma$ or $\gamma X$, e.g.\ with $X = Z$. In this case, the spectrum develops a prominent line centered at $E_{\gamma}=M_{DM}$.
    \item $\gamma$-ray continuum, from prompt emission (hadronization and decay of annihilation products) and from secondary radiation processes.
\end{itemize}

Supposing a signal in gamma ray lines will be observed, then the size of the continuum with respect to the
line and its spread in energy could allow to distinguish between different models in explaining the
origin of the line, if the relative background can
be understood.
Through observations of $\gamma$-rays (high energy photons), it is possible to derive some of the most stringent constraints on the DM annihilation cross section. 
The main advantage of gamma rays is the fact that photons are not deflected by magnetic fields and are negligibly attenuated over Galactic distance scales.
This allows to retain both spatial and spectral information, largely unaffected by astrophysical effects.
Assuming self-conjugate DM, the differential $\gamma$-ray flux per unit solid angle and energy from DM annihilation, integrated over the line of sight $s$, can be written as

\begin{equation}
\label{eq:AnnihilationFlux_general}
\frac{d\Phi_\gamma}{d\Omega\, dE}
=
\frac{1}{4\pi}
\frac{1}{2 M_{\mathrm{DM}}^2}
\sum_f \langle\sigma v\rangle_f\,
\frac{dN_{f\to\gamma}}{dE}
\int_{\text{l.o.s.}} ds\;
\rho_{\mathrm{DM}}^2(\vec r)\, .
\end{equation}

This expression, derived in appendix~\ref{app:gammaflux}, connects the microphysical properties of Dark Matter to the astrophysical factors, which depend solely on the DM density profile along the line of sight.
The quantity described by the integral in \cref{eq:AnnihilationFlux_general} is often referred to as the $J$ factor, which encodes all of the relevant astrophysical information.
In the particular case of the Milky Way it is convenient to define a dimensionless $J$ factor, normalized over the distance from the Earth to the Galactic Center $r_\odot$ and the local DM density $\rho_\odot$

\begin{equation} \label{Jfactor}
    J_{ann} (\theta) = \int_{l.o.s} \frac{ds}{r_\odot} \left( \frac{\rho_{DM}(r(s, \theta))}{\rho_\odot} \right)^2.
\end{equation}

The radial coordinate $r$ is measured from the GC and is related to $s$ through $r = \sqrt{r_\odot^{2} + s^{2} - 2\,r_\odot\, s \cos\theta}\,$. $\theta$ is the aperture angle between the direction of the line of sight and the axis linking the Earth to the GC.

The $J$ factors for different DM profiles are plotted in \cref{Jfactor_plot} as a function of $\theta$.


\begin{figure}[H]
\centering
\includegraphics[scale=0.50]{J-factor.png}
\caption{\emph{$J$ factors for different DM profiles. From top to bottom, the color individuates the profiles: Moore \cite{Moore1999}, NFW \cite{NFW1996,NFW1997}, EinastoB, Einasto \cite{Einasto1965}, Iso \cite{Begeman1991}, Burkert \cite{Burkert1995}. 
               Einasto and EinastoB have the same functional form but different value of the parameter \cite{Marco_Cirelli_2011}. } }
\label{Jfactor_plot}
\end{figure}





Eq.~\ref{Jfactor} is useful if one needs the flux of gamma rays from a given direction. Often one needs instead the flux over solid angle regions, for example, an annulus centered around the GC with $\theta_{\min} < \theta < \theta_{\max}$

\begin{equation}
\Delta\Omega
= 2\pi \int_{\theta_{\min}}^{\theta_{\max}} d\theta\,\sin\theta,
\qquad
\bar{J}
= \frac{2\pi}{\Delta\Omega}\int_{\theta_{\min}}^{\theta_{\max}} d\theta\,\sin\theta\, J(\theta).
\label{eq:DeltaOmega_Jbar}
\end{equation}

From the definition of the $J$ factor, one can infer the main features of ideal targets for Dark Matter detection: they should exhibit a high DM density, be located at relatively small distances, be spatially extended and feature a low or well understood astrophysical background \cite{hooper2018tasi}.
Dwarf spheroidal galaxies are particularly interesting because, although their $J$ factor is typically smaller than that of the inner volume of the Milky Way (the Galactic Center), they are characterized by a negligible astrophysical background. Nevertheless, they present some observational limitations: the so-called \emph{ultra-faint dwarfs} allow spectroscopic analyses for only a limited number of stars and the low $J$ factor implies that years of studies are required to obtain sufficient data to draw meaningful conclusions.
On the other hand, the Galactic Center is certainly the most promising source for Dark Matter annihilation in the Milky Way, however it presents a large and not fully understood astrophysical background.


\subsection{Galactic Center}

The Galactic Center is expected to be the single brightest source of dark matter annihilation products on the sky. In particular, using a standard NFW profile, shown in \cref{NFW}, with $r_s = 20$
kpc and a local density of $0.4~\text{GeV/cm}^3$, from \cref{eq:AnnihilationFlux_general} one obtains a flux of gamma rays that is more than three orders of magnitude larger than that predicted
from the brightest dwarf galaxies.
However, this mass whithin the innermost kpc of the Milky Way is dominated by baryons, 
this leads to an higher astrophysical background and, most importantly, numerical simulations need to include hydrodynamics
and feedback physics in addition to the gravitational effects, resulting in uncertainties in the obtained DM profiles towards the GC.
These simulations agree that the DM profile of MW-like galaxies develops a peak towards the Galactic Center but their resolution prevents them from determining a well-defined Dark Matter distribution.

As long as a clearer determination of the DM profile in the inner MW regions is not achieved, it is worth considering a range of possible choices and compare them. 
In particular, the two limiting cases are: a peaked profile, in which the density diverges as a power law towards the GC, and a cored one, in which the density flattens to a finite central value. 
Therefore, in what follows it is considered the Einasto (peaked) \cite{Einasto1965}, Burkert (cored) \cite{Burkert1995} 
and NFW \cite{NFW1996,NFW1997} with a variable core radius $r_c$ as representative DM density profiles 
$\rho_\mathrm{DM}(r)$.


For the Einasto and Burkert profiles, the functional forms read

\begin{equation}
\label{eq:rho_einasto}
\rho_{\mathrm{Ein}}(r)
=
\rho_s \exp\!\left[
-\frac{2}{\alpha}\left(\left(\frac{r}{r_s}\right)^{\alpha}-1\right)
\right].
\end{equation}

\begin{equation}
\label{eq:rho_burkert}
\rho_{\mathrm{Bur}}(r)
=
\frac{\rho_s}{\left(1+\dfrac{r}{r_s}\right)\left(1+\left(\dfrac{r}{r_s}\right)^2\right)}.
\end{equation}

where $\alpha = 0.17$.

The functional form of the cored NFW is


\begin{equation}
\label{eq:rho_nfw_cored}
\rho_{\mathrm{NFW}}(r,r_c)=
\begin{cases}
\rho_s\,\dfrac{r_s}{r_c}\,\dfrac{1}{\left(1+\dfrac{r_c}{r_s}\right)^2},
& \text{for } r \le r_c, \\[10pt]
\rho_s\,\dfrac{r_s}{r}\,\dfrac{1}{\left(1+\dfrac{r}{r_s}\right)^2},
& \text{for } r > r_c.
\end{cases}
\end{equation}

Once $r_c$ is fixed, all three profiles depend on two free parameters, $\rho_s$ and $r_s$.
They are determined by fitting the DM density at the Sun's position: $\rho(r_\odot = 8.33~\text{kpc}) = 0.3~\text{GeV/cm}^3$, and the total DM mass within 60~kpc: $M(<60~\text{kpc}) = 4.7 \times 10^{11}\,M_\odot$.

The $J$-factors are integrated over an angular region, centred on the GC, defined as 

\begin{equation}
\label{eq:Omega_obs}
    \Omega_\mathrm{obs} = \Omega(\theta < 1^\circ,\; |b| > 0.3^\circ),
\end{equation}

where $\theta > 0^\circ$ is the polar angle with respect to the GC direction 
and $b$ is the galactic latitude with respect to the disk plane. At the Sun's 
distance $r_\odot = 8.33$~kpc, the angles $1^\circ$ and $0.4^\circ$ correspond 
to projected distances of approximately 145~pc and 44~pc from the GC, respectively. 
The choice of $|b| > 0.3^\circ$ is motivated by the need to mask the large photon backgrounds from the galactic ridge.
The results are reported in \cref{tab:gc_jfactors}, where for the cored NFW has been used three illustrative values of $r_c$.


\begin{table}[H]
\centering
\begin{tabular}{c|c|c|c|c|c}
\hline
 & Einasto & Burkert & $\mathrm{NFW}_{30\,\mathrm{pc}}$ & $\mathrm{NFW}_{300\,\mathrm{pc}}$ & $\mathrm{NFW}_{3\,\mathrm{kpc}}$ \\
\hline
$\log_{10}\!\left(J_{\Omega_{\mathrm{obs}}}\,\left[\dfrac{\mathrm{GeV}^2}{\mathrm{cm}^5}\right]\right)$
& 21.96 & 19.35 & 21.57 & 21.16 & 19.77 \\
\hline
\end{tabular}
\caption{Galactic Center integrated $J$-factors, with NFW core radii reported as pedices \cite{Lefranc_2016}.}
\label{tab:gc_jfactors}
\end{table}

The values in \cref{tab:gc_jfactors} span roughly two orders of magnitude 
between the Einasto, the most peaked profile, and Burkert, the most cored one.
This has a direct consequence for 
the upper bounds on the DM annihilation cross-section: since the predicted flux scales linearly with the 
$J$-factor, the upper bound on $\langle \sigma v \rangle$ is inversely 
proportional to it. Consequently, a more peaked profile yields a stronger constraint than a cored one.
In particular, the upper bound on $\langle \sigma v \rangle$ derived under the Burkert assumption 
is roughly a factor $\sim 4 \times 10^2$ weaker than the one obtained under Einasto.



